% Options for packages loaded elsewhere
\PassOptionsToPackage{unicode}{hyperref}
\PassOptionsToPackage{hyphens}{url}
%
\documentclass[
]{article}
\usepackage{amsmath,amssymb}
\usepackage[margin=1in]{geometry}
\usepackage{iftex}
\ifPDFTeX
  \usepackage[T1]{fontenc}
  \usepackage[utf8]{inputenc}
  \usepackage{textcomp} % provide euro and other symbols
\else % if luatex or xetex
  \usepackage{unicode-math} % this also loads fontspec
  \defaultfontfeatures{Scale=MatchLowercase}
  \defaultfontfeatures[\rmfamily]{Ligatures=TeX,Scale=1}
\fi
\usepackage{lmodern}
\ifPDFTeX\else
  % xetex/luatex font selection
\fi
% Use upquote if available, for straight quotes in verbatim environments
\IfFileExists{upquote.sty}{\usepackage{upquote}}{}
\IfFileExists{microtype.sty}{% use microtype if available
  \usepackage[]{microtype}
  \UseMicrotypeSet[protrusion]{basicmath} % disable protrusion for tt fonts
}{}
\makeatletter
\@ifundefined{KOMAClassName}{% if non-KOMA class
  \IfFileExists{parskip.sty}{%
    \usepackage{parskip}
  }{% else
    \setlength{\parindent}{0pt}
    \setlength{\parskip}{6pt plus 2pt minus 1pt}}
}{% if KOMA class
  \KOMAoptions{parskip=half}}
\makeatother
\usepackage{xcolor}
\usepackage{color}
\usepackage{fancyvrb}
\newcommand{\VerbBar}{|}
\newcommand{\VERB}{\Verb[commandchars=\\\{\}]}
\DefineVerbatimEnvironment{Highlighting}{Verbatim}{commandchars=\\\{\}}
% Add ',fontsize=\small' for more characters per line
\newenvironment{Shaded}{}{}
\newcommand{\AlertTok}[1]{\textcolor[rgb]{1.00,0.00,0.00}{\textbf{#1}}}
\newcommand{\AnnotationTok}[1]{\textcolor[rgb]{0.38,0.63,0.69}{\textbf{\textit{#1}}}}
\newcommand{\AttributeTok}[1]{\textcolor[rgb]{0.49,0.56,0.16}{#1}}
\newcommand{\BaseNTok}[1]{\textcolor[rgb]{0.25,0.63,0.44}{#1}}
\newcommand{\BuiltInTok}[1]{\textcolor[rgb]{0.00,0.50,0.00}{#1}}
\newcommand{\CharTok}[1]{\textcolor[rgb]{0.25,0.44,0.63}{#1}}
\newcommand{\CommentTok}[1]{\textcolor[rgb]{0.38,0.63,0.69}{\textit{#1}}}
\newcommand{\CommentVarTok}[1]{\textcolor[rgb]{0.38,0.63,0.69}{\textbf{\textit{#1}}}}
\newcommand{\ConstantTok}[1]{\textcolor[rgb]{0.53,0.00,0.00}{#1}}
\newcommand{\ControlFlowTok}[1]{\textcolor[rgb]{0.00,0.44,0.13}{\textbf{#1}}}
\newcommand{\DataTypeTok}[1]{\textcolor[rgb]{0.56,0.13,0.00}{#1}}
\newcommand{\DecValTok}[1]{\textcolor[rgb]{0.25,0.63,0.44}{#1}}
\newcommand{\DocumentationTok}[1]{\textcolor[rgb]{0.73,0.13,0.13}{\textit{#1}}}
\newcommand{\ErrorTok}[1]{\textcolor[rgb]{1.00,0.00,0.00}{\textbf{#1}}}
\newcommand{\ExtensionTok}[1]{#1}
\newcommand{\FloatTok}[1]{\textcolor[rgb]{0.25,0.63,0.44}{#1}}
\newcommand{\FunctionTok}[1]{\textcolor[rgb]{0.02,0.16,0.49}{#1}}
\newcommand{\ImportTok}[1]{\textcolor[rgb]{0.00,0.50,0.00}{\textbf{#1}}}
\newcommand{\InformationTok}[1]{\textcolor[rgb]{0.38,0.63,0.69}{\textbf{\textit{#1}}}}
\newcommand{\KeywordTok}[1]{\textcolor[rgb]{0.00,0.44,0.13}{\textbf{#1}}}
\newcommand{\NormalTok}[1]{#1}
\newcommand{\OperatorTok}[1]{\textcolor[rgb]{0.40,0.40,0.40}{#1}}
\newcommand{\OtherTok}[1]{\textcolor[rgb]{0.00,0.44,0.13}{#1}}
\newcommand{\PreprocessorTok}[1]{\textcolor[rgb]{0.74,0.48,0.00}{#1}}
\newcommand{\RegionMarkerTok}[1]{#1}
\newcommand{\SpecialCharTok}[1]{\textcolor[rgb]{0.25,0.44,0.63}{#1}}
\newcommand{\SpecialStringTok}[1]{\textcolor[rgb]{0.73,0.40,0.53}{#1}}
\newcommand{\StringTok}[1]{\textcolor[rgb]{0.25,0.44,0.63}{#1}}
\newcommand{\VariableTok}[1]{\textcolor[rgb]{0.10,0.09,0.49}{#1}}
\newcommand{\VerbatimStringTok}[1]{\textcolor[rgb]{0.25,0.44,0.63}{#1}}
\newcommand{\WarningTok}[1]{\textcolor[rgb]{0.38,0.63,0.69}{\textbf{\textit{#1}}}}
\usepackage{longtable,booktabs,array}
\usepackage{calc} % for calculating minipage widths
% Correct order of tables after \paragraph or \subparagraph
\usepackage{etoolbox}
\makeatletter
\patchcmd\longtable{\par}{\if@noskipsec\mbox{}\fi\par}{}{}
\makeatother
% Allow footnotes in longtable head/foot
\IfFileExists{footnotehyper.sty}{\usepackage{footnotehyper}}{\usepackage{footnote}}
\makesavenoteenv{longtable}
\setlength{\emergencystretch}{3em} % prevent overfull lines
\providecommand{\tightlist}{%
  \setlength{\itemsep}{0pt}\setlength{\parskip}{0pt}}
\setcounter{secnumdepth}{-\maxdimen} % remove section numbering
\ifLuaTeX
  \usepackage{selnolig}  % disable illegal ligatures
\fi
\usepackage{bookmark}
\IfFileExists{xurl.sty}{\usepackage{xurl}}{} % add URL line breaks if available
\urlstyle{same}
\hypersetup{
  hidelinks,
  pdfcreator={LaTeX via pandoc}}

\author{}
\date{}

\title{The General Divisor Theorem: Exact Density Correction under Residue Conditioning}
\author{thinkthoughts}
\date{}
\begin{document}
\maketitle

\section{The General Divisor Theorem: Exact Density Correction under
Residue
Conditioning}\label{the-general-divisor-theorem-exact-density-correction-under-residue-conditioning}

thinkthoughts https://github.com/thinkthoughts/mathematical-basis

\textbf{Draft v3.2} --- extends Draft v3.1 (preserved unchanged as
\texttt{paper-v3.1.pdf}); see the Change Log at the end of this document
for a full record of changes.

\subsection{Abstract}\label{abstract}

Let \(N,m\) be positive integers and \(a\) an integer. Write \[
d=\mathrm{rad}(\gcd(m,N)),
\qquad
R=\frac{\mathrm{rad}(N)}{d}.
\] For \[
S(N,L,m,a)
=
\{1\le n\le L:n\equiv a\pmod m,\ \gcd(n,N)=1\},
\] we prove an exact dichotomy. If \(\gcd(a,d)=1\), the associated
indicator has exact minimal period \[
T_{\min}=mR=\mathrm{lcm}(m,\mathrm{rad}(N)),
\] with exactly \(\varphi(R)\) accepted values per period. Relative to
the naive independent-density prediction \[
\frac{\varphi(N)}{N}\frac{L}{m},
\] the exact correction factor is \[
C(N,m)=\frac{d}{\varphi(d)},
\] independent of the admissible residue \(a\). If \(\gcd(a,d)>1\), the
accepted set is empty for every \(L\). These period, count, and
correction formulas are exact, rather than asymptotic, and require no
divisibility or squarefreeness restriction on \(N\) or \(m\).

The theorem recovers the motivating primorial result when \(N=P_k\) and
\(m\mid P_k\): squarefreeness gives \(d=m\), \(R=P_k/m\),
\(T_{\min}=P_k\), and \(C(N,m)=m/\varphi(m)\). We also examine a
separate weighted-vector construction as a worked example distinguishing
a shared label, a shared reading, and a shared specification.

\subsection{1. Introduction / General
Setup}\label{introduction-general-setup}

Let \(N\) and \(m\) be positive integers and \(a\) an integer. We study
the counting problem \[
S(N,L,m,a) = \{1\le n\le L : n\equiv a\pmod m,\ \gcd(n,N)=1\},
\] the integers up to \(L\) lying in a fixed residue class modulo \(m\)
and coprime to \(N\).

A naive density heuristic treats the residue condition and the
coprimality condition as independent: density \(1/m\) from the residue
class, multiplied by \(\varphi(N)/N\) from coprimality to \(N\). This
gives the naive prediction \[
|S(N,L,m,a)| \approx \frac{\varphi(N)}{N}\frac{L}{m}.
\]

This double-counts a restriction whenever \(m\) and \(N\) share a prime
factor. If a prime \(p\) divides both \(m\) and \(N\), and
\(n\equiv a\pmod m\) with \(\gcd(a,p)=1\), then \(n\) is already
guaranteed coprime to \(p\) by the residue condition alone --- the
factor \(\varphi(N)/N\) re-imposes that restriction a second time. The
primes doing this double duty are exactly the primes shared between
\(m\) and \(N\); primes of \(N\) not dividing \(m\), and primes of \(m\)
not dividing \(N\), play no special role in this double-counting.

This motivates isolating the shared-prime structure directly, via \[
d=\mathrm{rad}(\gcd(m,N)), \qquad R=\frac{\mathrm{rad}(N)}{d},
\] the radical of the primes \(m\) and \(N\) have in common, and the
radical formed from the primes of \(N\) not shared with \(m\). Section 2
states and proves the exact correction this produces. Section 3 walks
through its specializations, recovering progressively more restrictive
--- and historically earlier --- versions of the same result, down to
the primorial case that originally motivated this work.

\subsection{2. The General Divisor
Theorem}\label{the-general-divisor-theorem}

\subsubsection{2.1 A lemma used
throughout}\label{a-lemma-used-throughout}

\textbf{Lemma} (exponent-blindness). For every integer \(k\ge1\), \[
\frac{k}{\varphi(k)} = \frac{\mathrm{rad}(k)}{\varphi(\mathrm{rad}(k))},
\] where \(\mathrm{rad}(k)\) is the product of the distinct primes
dividing \(k\) (\(\mathrm{rad}(1)=1\)).

\textbf{Proof.} For \(k=1\) both sides equal \(1\). For \(k>1\), write
\(k=\prod_i p_i^{e_i}\). For each prime power,
\(\varphi(p_i^{e_i})/p_i^{e_i}=1-1/p_i\), independent of \(e_i\). By
multiplicativity, \[
\frac{k}{\varphi(k)} = \prod_i\frac{1}{1-1/p_i} = \prod_{p\mid k}\frac{p}{p-1},
\] which depends only on the set of primes dividing \(k\) --- i.e.~only
on \(\mathrm{rad}(k)\). Applying the identical computation to
\(\mathrm{rad}(k)\) (all exponents \(1\)) gives the same product.
\(\blacksquare\)

\subsubsection{2.2 Theorem 3 (General Divisor
Theorem)}\label{theorem-3-general-divisor-theorem}

Let \(N,m\) be any positive integers, \(a\) any integer. Define \[
d = \mathrm{rad}(\gcd(m,N)), \qquad R = \frac{\mathrm{rad}(N)}{d}.
\] Let \(S(L)=\{1\le n\le L : n\equiv a\pmod m,\ \gcd(n,N)=1\}\).

\textbf{Case \(\gcd(a,d)=1\).} Writing \(L=qT_{\min}+s\),
\(0\le s<T_{\min}\), \[
|S(L)| = q\,\varphi(R) + R_{\mathrm{rem}}(s), \qquad T_{\min} = mR = \mathrm{lcm}(m,\mathrm{rad}(N)),
\] \(T_{\min}\) is the exact minimal period of the indicator
\(\mathbb1[n\equiv a\pmod m,\ \gcd(n,N)=1]\), and against the naive
predicted count \(\frac{\varphi(N)}{N}\frac{L}{m}\), \[
C(N,m) = \frac{d}{\varphi(d)} = \prod_{p\mid\gcd(m,N)}\frac{p}{p-1}.
\]

\textbf{Case \(\gcd(a,d)>1\).} Then \(S(L)=\varnothing\) for every
\(L\).

\textbf{Proof.}

\emph{Dichotomy and exact count.} Fix \(n\equiv a\pmod m\). For each
prime \(p\mid d\): since \(p\mid m\), \(n\equiv a\pmod p\), so \(n\) is
coprime to \(p\) iff \(a\) is. If \(\gcd(a,d)=1\), every such \(n\) is
coprime to all of \(d\); if some \(p\mid d\) divides \(a\), then \(p\)
divides every \(n\) in the class, and since \(p\mid d\mid N\), no such
\(n\) is coprime to \(N\) --- giving \(S(L)=\varnothing\) for every
\(L\). Assume now \(\gcd(a,d)=1\): then \(\gcd(n,N)=1 \iff \gcd(n,R)=1\)
(the only primes of \(N\) left to check are \(R\)'s). Since
\(\gcd(m,R)=1\), any interval of \(mR\) consecutive integers is a
complete residue system mod \(mR\); among such an interval, the
sub-collection with \(n\equiv a\pmod m\) numbers exactly \(R\) elements
and bijects with \(\mathbb Z/R\mathbb Z\) via \(n\bmod R\) (CRT), of
which exactly \(\varphi(R)\) are coprime to \(R\). This holds for every
such window, giving the exact count \(\varphi(R)\) per
length-\(T_{\min}\) window and, summing over \(q\) disjoint windows,
\(|S(qT_{\min})|=q\varphi(R)\).

\emph{\(T_{\min}\) is a period.} For \(n\equiv a\pmod m\),
\(\gcd(n,N)=1\iff\gcd(n,R)=1\) as above; for \(n\not\equiv a\pmod m\),
\(n\) is unaccepted regardless. So the indicator factors as
\(I(n)=\mathbb1[n\equiv a\pmod m]\cdot\mathbb1[\gcd(n,R)=1]\) for all
\(n\). The first factor has period \(m\) (depends only on \(n\bmod m\));
the second has period \(R\) (coprimality to \(R\) depends only on
\(n\bmod R\), since \(R\) is squarefree). Since \(m\mid mR\) and
\(R\mid mR\), both factors, and hence their product, are unchanged under
\(n\mapsto n+mR\): \(T_{\min}=mR\) is a period.

\emph{No smaller \(t\) is a period.} By CRT (using \(\gcd(m,R)=1\)) pick
\(n_0\equiv a\pmod m\), \(n_0\equiv1\pmod R\); \(n_0\) is accepted. If
\(t>0\) is any period, \(n_0+t\) is accepted, so
\(n_0+t\equiv a\pmod m\), giving \(m\mid t\). Fix a prime \(p\mid R\)
and suppose \(p\nmid t\); by CRT choose \(n_1\equiv a\pmod m\),
\(n_1\equiv-t\pmod p\), \(n_1\equiv1\pmod q\) for every other prime
\(q\mid R\). Then \(n_1\) is accepted but \(n_1+t\equiv0\pmod p\) is not
--- contradicting periodicity. Hence \(p\mid t\) for every prime
\(p\mid R\); since \(R\) is squarefree, \(R\mid t\), and with
\(m\mid t\), \(\gcd(m,R)=1\): \(mR\mid t\). Thus \(mR\mid t\) for every
positive period \(t\). Since \(mR\) itself is a period, it is the exact
minimal positive period: \(T_{\min}=mR\).

\emph{Ratio.} Write \(N=AB\), \(A=\prod_{p\mid N,p\mid m}p^{e_p}\)
(\(\mathrm{rad}(A)=d\)), \(B=\prod_{p\mid N,p\nmid m}p^{e_p}\)
(\(\mathrm{rad}(B)=R\)), \(\gcd(A,B)=1\), so
\(\varphi(N)=\varphi(A)\varphi(B)\). By the Lemma (applied to \(k=B\)),
\(\varphi(B)/B=\varphi(R)/R\). Then \[
C = \frac{\varphi(R)/(mR)}{\varphi(N)/(Nm)} = \frac{N\varphi(R)}{R\varphi(N)} = \frac{AB\varphi(R)}{R\varphi(A)\varphi(B)} = \frac{AB\varphi(R)}{R\varphi(A)(B\varphi(R)/R)} = \frac{A}{\varphi(A)}.
\] By the Lemma once more (applied to \(k=A\)),
\(A/\varphi(A)=d/\varphi(d)\), giving \(C=d/\varphi(d)\).
\(\blacksquare\)

\subsection{3. Specializations and Distinguishing
Examples}\label{specializations-and-distinguishing-examples}

Having proved Theorem 3 in full generality, we now walk down its
specializations, each obtained by substituting a stronger hypothesis on
\((N,m,a)\) directly into \(d\), \(R\), \(T_{\min}\), and \(C(N,m)\) ---
no independent re-derivation is required at any stage: \[
\text{Theorem 3} \;\longrightarrow\; \text{Theorem 2} \;\longrightarrow\; \text{Theorem }1' \;\longrightarrow\; \text{Theorem 1},
\] where each arrow denotes ``specializes to'': the target theorem
imposes stronger hypotheses, and its period, count, and
correction-factor formulas are recovered from the source theorem by
direct substitution, not by an independent proof.

\subsubsection{\texorpdfstring{3.1 Theorem 2 (specialization:
\(m\mid N\))}{3.1 Theorem 2 (specialization: m\textbackslash mid N)}}\label{theorem-2-specialization-mmid-n}

If \(m\mid N\), then \(\gcd(m,N)=m\), so \(d=\mathrm{rad}(m)\) directly.
Substituting into Theorem 3: \[
R = \frac{\mathrm{rad}(N)}{\mathrm{rad}(m)}, \qquad T_{\min} = m\,\frac{\mathrm{rad}(N)}{\mathrm{rad}(m)}, \qquad C(N,m) = \frac{\mathrm{rad}(m)}{\varphi(\mathrm{rad}(m))} = \frac{m}{\varphi(m)}
\] (the last equality by the Lemma applied to \(k=m\)), with
admissibility \(\gcd(a,\mathrm{rad}(m))=1\) and the same complementary
zero branch. No new argument is required beyond this substitution;
minimality, the exact count \(\varphi(R)\), and the ratio are all
inherited directly from Theorem 3.

\subsubsection{\texorpdfstring{3.2 Theorem~\(1'\) (specialization: \(m\)
unitary in
\(N\))}{3.2 Theorem 1' (specialization: m unitary in N)}}\label{theorem-1-specialization-m-unitary-in-n}

Call \(m\) a \emph{unitary divisor} of \(N\) if \(\gcd(m,N/m)=1\) ---
equivalently, \(m\) absorbs the full exponent of each of its primes in
\(N\). Under this hypothesis (a strengthening of \(m\mid N\)), write
\(N'=N/m\). Then \(\mathrm{rad}(N)=\mathrm{rad}(m)\mathrm{rad}(N')\)
(coprime factors), so
\(R=\mathrm{rad}(N)/\mathrm{rad}(m)=\mathrm{rad}(N')\) and \[
T_{\min} = m\mathrm{rad}(N'), \qquad C(N,m)=\frac{m}{\varphi(m)},
\] as in Theorem 2. The interval \(N=mN'\) is also a period of the
indicator (it is an integer multiple of \(T_{\min}\), since
\(N'/\mathrm{rad}(N')\) is an integer), with total count \(\varphi(N')\)
over that longer interval --- this matches Theorem 3's prediction
exactly (using exponent-blindness to check
\(\varphi(N')/N'=\varphi(\mathrm{rad}(N'))/\mathrm{rad}(N')\)), but
\textbf{\(N\) is generally not the minimal period}: it equals
\(T_{\min}\) only when \(N'\) is itself squarefree.

\subsubsection{3.3 Two distinguishing
examples}\label{two-distinguishing-examples}

\textbf{\(N=18,\ m=2\) --- coarser vs.~minimal period.} Here \(m=2\) is
a unitary divisor of \(N=18=2\cdot3^2\) (\(\gcd(2,9)=1\)), so Theorem~\(1'\)
applies and gives a valid period of \(N=18\) with count
\(\varphi(9)=6\). But \(\mathrm{rad}(N')=\mathrm{rad}(9)=3\), so Theorem
3's exact minimal period is \[
T_{\min} = m\mathrm{rad}(N') = 2\cdot3 = 6.
\] Direct enumeration confirms this: taking \(a=1\), \(n\equiv1\pmod2\)
and \(\gcd(n,18)=1\) (odd and not divisible by \(3\)) gives
\(\{1,5,7,11,13,17,\dots\}\) in \([1,18]\), i.e.~\(\{1,5,7,11,13,17\}\)
--- a pattern that visibly repeats every \(6\) integers, not every
\(18\). So \(18\) is a genuine period (as Theorem~\(1'\) correctly states)
but not the minimal one; \(6\) is.

\textbf{\(N=5,\ m=6,\ a=2\) --- the sharp admissibility condition.} Here
\(\gcd(a,m)=\gcd(2,6)=2\ne1\), so this case would be \emph{excluded}
under Theorem 1's or Theorem~\(1'\)'s hypothesis \(\gcd(a,m)=1\). But
\(d=\mathrm{rad}(\gcd(6,5))=\mathrm{rad}(1)=1\), and
\(\gcd(a,d)=\gcd(2,1)=1\): admissible under Theorem 3's sharp condition.
Here \(R=\mathrm{rad}(5)=5\) (since \(5\nmid6\)), \(T_{\min}=mR=30\).
Direct check, \(n\equiv2\pmod6\) in \([1,30]\): \(\{2,8,14,20,26\}\);
coprime to \(5\): \(\{2,8,14,26\}\) --- count \(4=\varphi(5)\), matching
the theorem. And \(C(5,6)=d/\varphi(d)=1\), matching the naive
prediction exactly (\(\varphi(5)/5\cdot(30/6)=4\)) --- whereas
\(m/\varphi(m)=6/\varphi(6)=3\) would have been wrong. This confirms
that \(\gcd(a,d)=1\) is the exact admissibility condition, not merely a
cosmetic restatement of \(\gcd(a,m)=1\): it is strictly weaker in
general, admitting strictly more values of \(a\), since \(d\) contains
only the primes shared between \(m\) and \(N\).

\subsubsection{\texorpdfstring{3.4 Theorem 1 (specialization: primorial
\(N\))}{3.4 Theorem 1 (specialization: primorial N)}}\label{theorem-1-specialization-primorial-n}

Let \(P_k = p_1p_2\cdots p_k\) denote the \(k\)-th primorial, the
product of the first \(k\) primes (\(p_1=2,p_2=3,\dots\)). By
construction \(P_k\) is squarefree. Let \(m\) be a squarefree divisor of
\(P_k\) with \(m>1\), and let \(a\) be an integer with \(\gcd(a,m)=1\).

If, in addition to \(m\mid N\), \(N\) is squarefree (\(N=P_k\)), then
every divisor of \(N\) --- including \(m\) --- is automatically unitary,
so Theorem~\(1'\) applies with \(N'=P_k/m\); and since \(P_k\) is
squarefree, \(N'=P_k/m\) is squarefree too (a divisor of a squarefree
number is squarefree), so \(\mathrm{rad}(N')=N'\) and the coarser and
minimal periods of Section 3.2 coincide exactly: \[
T_{\min}=m\mathrm{rad}(N')=mN'=P_k.
\] Because \(P_k\) is squarefree, \(\mathrm{rad}(m)=m\), so
\(d=\mathrm{rad}(\gcd(m,P_k))=\mathrm{rad}(m)=m\) exactly --- the sharp
admissibility condition \(\gcd(a,d)=1\) used in Theorem 3 is, in this
primorial setting, \emph{identical} to the hypothesis \(\gcd(a,m)=1\)
above. This is the same pattern already seen in Theorem 2
(\(d=\mathrm{rad}(m)\) whenever \(m\mid N\)); it collapses to an
equality here specifically because \(m\) is squarefree.

This gives:

\textbf{Theorem 1} (primorial specialization). Let \(P_k\), \(m\), and
\(a\) be as above. Write \[
L = qP_k + s,\qquad 0\le s<P_k.
\] Then \[
|S(P_k,L,m,a)| = q\,\varphi(P_k/m) + R(s),
\] where \[
R(s) = \bigl|\{1\le n\le s : n\equiv a\pmod m,\ \gcd(n,P_k)=1\}\bigr|
\] is the exact count within the partial final period and satisfies
\(0\le R(s)\le\lceil s/m\rceil\), with minimal period \(P_k\) and
correction factor \[
C(m) = \frac{m}{\varphi(m)},
\] independent of the admissible residue \(a\). This is exactly Theorem
3 specialized as above, with no adjustment to period, count, or
correction factor required.

\paragraph{Independent primorial
proof}\label{independent-primorial-proof}

The following is the original proof of Theorem 1, self-contained and
independent of Theorem 3, retained here as the paper's historical
starting point.

\textbf{Proof.} Since \(m\mid P_k\) and \(P_k\) is squarefree, \(m\) and
\(P_k/m\) are coprime and squarefree, and \(P_k/m\) consists exactly of
the prime factors of \(P_k\) that do not divide \(m\).

Within one period \([1,P_k]\), exactly \(P_k/m\) integers satisfy
\(n\equiv a\pmod m\). By the Chinese Remainder Theorem, because \(m\)
and \(P_k/m\) are coprime, the residue of \(n\) modulo \(P_k/m\) ranges
uniformly over all residue classes as \(n\) ranges over its fixed class
modulo \(m\) within one period.

Hence exactly the fraction \[
\prod_{p\mid P_k/m}\Bigl(1-\frac1p\Bigr) = \frac{\varphi(P_k/m)}{P_k/m}
\] of these \(P_k/m\) values are additionally coprime to every prime
dividing \(P_k/m\). Since \(\gcd(a,m)=1\), they are already coprime to
every prime dividing \(m\). Therefore the exact count per period is \[
\frac{P_k}{m}\cdot\frac{\varphi(P_k/m)}{P_k/m} = \varphi(P_k/m).
\]

The count does not depend on the specific value of \(a\) beyond
\(\gcd(a,m)=1\), because the argument uses only that the residue class
modulo \(m\) is fixed and coprime to \(m\).

Using multiplicativity of Euler's totient function,
\(\varphi(P_k)=\varphi(m)\varphi(P_k/m)\), so
\(\varphi(P_k/m)=\varphi(P_k)/\varphi(m)\). The naive predicted count
per period is \(\varphi(P_k)/m\). The ratio is therefore \[
\frac{\varphi(P_k)/\varphi(m)}{\varphi(P_k)/m} = \frac{m}{\varphi(m)} = C(m).
\] This ratio is exact per full period and independent of \(k\).
Periodicity with period \(P_k\) gives the general-\(L\) statement with
the bounded remainder above. \(\blacksquare\)

\paragraph{Standardized refined-density
identity}\label{standardized-refined-density-identity}

Throughout this paper, the refined predicted count is written in the
equivalent forms \[
E_{\mathrm{refined}}(L) = \frac{L}{m}\prod_{p\mid P_k/m}\Bigl(1-\frac1p\Bigr) = L\,\frac{\varphi(P_k/m)}{P_k}.
\] These are equal because
\(\frac{L}{m}\cdot\frac{\varphi(P_k/m)}{P_k/m} = L\,\frac{\varphi(P_k/m)}{P_k}\).

\textbf{Corollary 1} (\(m=6\) case). For \(m=6\) and \(a\in\{1,5\}\), \[
C(6) = \frac{6}{\varphi(6)} = \frac62 = 3.
\] This corrects the earlier draft's reported limiting constant
\(24/25\) for the case \(a=5\).

Against the refined predicted count
\(E_{\mathrm{refined}}(L)=L\varphi(P_k/6)/P_k\), the ratio is exactly
\(1\).

\emph{(In Section 2's general notation,
\(C(6)=6/\varphi(6)=d/\varphi(d)\) with \(d=6\), since \(6\) is already
squarefree --- the two notations agree exactly, as they must by the
remark above.)}

\textbf{Corollary 2} (Combined residue classes). Since the per-period
count \(\varphi(P_k/m)\) is independent of which admissible \(a\) is
used, the same ratio \(C(m)\) holds for any union of admissible residue
classes modulo \(m\). For example, the classes \(n\equiv1\) or
\(5\pmod6\) together give the same ratio \(3\) as either class
individually, because both the exact and naive counts scale by the same
number of residue classes.

Both corollaries were verified numerically at \(P_k=210\) for
\(m\in\{6,10,30\}\), multiple values of \(a\), and \(L\) up to \(10^7\),
matching the predicted \(C(m)\) to five decimal places or better in
every case and matching the exact per-period counts \(\varphi(P_k/m)\)
precisely; see \texttt{tests/basis/test\_persistent\_constant.py}.

\subsection{4. A Failed Correspondence as a Reading-Point
Example}\label{a-failed-correspondence-as-a-reading-point-example}

An earlier draft defined a weighted tuple \((9,4,2,3)\) with angles
\(\theta_k\in\{0^\circ,60^\circ,120^\circ,180^\circ\}\) and \[
V = \sum_{k=1}^4 w_k e^{i\theta_k} = 9 + 4e^{i\pi/3} + 2e^{i2\pi/3} + 3e^{i\pi} = 7 + 3\sqrt3\,i.
\] The reported reading was \(r_V = \arg(V) \approx 36.586776^\circ\).

The same draft separately noted that the point \((1,1)\) has reading
\(r_D=\arg(1+i)=45^\circ\), and described the combination as a ``9423
phase-lock representation,'' implying a correspondence between the two
constructions.

Since \(r_V\ne r_D\), no such correspondence holds under the stated
construction.

An exhaustive search over all \(24\) permutations of assigning the four
weights to the four fixed angles confirms that this discrepancy is not
an artifact of ordering: no permutation reaches exactly \(45^\circ\), with
the closest approaches at \(43.898^\circ\) and \(46.102^\circ\)
(\texttt{tests/basis/test\_9423\_phase\_lock.py}).

This gives the three-level distinction \[
\text{shared label} \not\Rightarrow \text{shared reading} \not\Rightarrow \text{shared specification.}
\]

The draft's own label, ``9423 phase-lock representation,'' was shared
across both objects by assertion, but the named readings are not shared:
\(r_V=36.586776^\circ \ne 45^\circ = r_D\).

The second implication is worth stating independently. Even had
\(r_V=r_D\) exactly, that equality alone would not establish a
structural correspondence between \(V\)'s construction and the point
\((1,1)\). No transformation, projection, or limiting operation
connecting the two objects would thereby be specified.

A shared label asserts a correspondence. A shared reading is, at most,
evidence that two objects were run through procedures producing the same
observable category. Neither, on its own, specifies the correspondence
itself.

Section 3 gives a cleaner instance of the same principle. The reported
ratio for \(m=6,a=5\) has taken three different numerical values during
the paper's development: \(24/25\) (originally reported, against no
explicitly stated refined predicted count), \(3\) (against the naive
predicted count \((\varphi(P_k)/P_k)(L/6)\)), and \(1\) (against the
refined predicted count \(L\varphi(P_k/6)/P_k\)). None of these numbers
is meaningful without specifying which predictor produced it. The number
alone is insufficient; the reading requires its producing specification
to be stated alongside it.

\subsection{5. Numerical Verification}\label{numerical-verification}

At \(P_k=210=2\cdot3\cdot5\cdot7\), \(m=6\), \(a=5\), the numerical
results are:

\textbf{Table 1: Numerical verification at \(P_k=210\), \(m=6\),
\(a=5\).}

\begin{longtable}[]{@{}
  >{\raggedright\arraybackslash}p{(\columnwidth - 10\tabcolsep) * \real{0.1667}}
  >{\raggedright\arraybackslash}p{(\columnwidth - 10\tabcolsep) * \real{0.1667}}
  >{\raggedright\arraybackslash}p{(\columnwidth - 10\tabcolsep) * \real{0.1667}}
  >{\raggedright\arraybackslash}p{(\columnwidth - 10\tabcolsep) * \real{0.1667}}
  >{\raggedright\arraybackslash}p{(\columnwidth - 10\tabcolsep) * \real{0.1667}}
  >{\raggedright\arraybackslash}p{(\columnwidth - 10\tabcolsep) * \real{0.1667}}@{}}
\toprule\noalign{}
\begin{minipage}[b]{\linewidth}\raggedright
\(L\)
\end{minipage} & \begin{minipage}[b]{\linewidth}\raggedright
actual
\end{minipage} & \begin{minipage}[b]{\linewidth}\raggedright
naive predicted
\end{minipage} & \begin{minipage}[b]{\linewidth}\raggedright
naive ratio
\end{minipage} & \begin{minipage}[b]{\linewidth}\raggedright
refined predicted
\end{minipage} & \begin{minipage}[b]{\linewidth}\raggedright
refined ratio
\end{minipage} \\
\midrule\noalign{}
\endhead
\bottomrule\noalign{}
\endlastfoot
\(10^5\) & 11,428 & 3,809.52 & 2.99985 & 11,428.57 & 0.99995 \\
\(10^6\) & 114,285 & 38,095.24 & 2.99998 & 114,285.71 & 0.99999 \\
\(10^7\) & 1,142,856 & 380,952.38 & 3.00000 & 1,142,857.14 & 1.00000 \\
\(10^8\) & 11,428,571 & 3,809,523.81 & 3.00000 & 11,428,571.43 &
1.00000 \\
\end{longtable}

Additional verification across \(m\in\{6,10,30\}\) and all admissible
\(a\bmod m\) confirms \(C(m)=m/\varphi(m)\) exactly: \(C(6)=3\),
\(C(10)=2.5\), \(C(30)=3.75\). It also confirms the exact per-period
count \(\varphi(P_k/m)\), independently of \(a\), at \(P_k=210\) and
\(L\) up to \(10^7\); see
\texttt{tests/basis/test\_persistent\_constant.py}.

\emph{(These values also follow directly from Section 2's general
formulas: e.g. \(C(10)=10/\varphi(10)=2.5\) and
\(C(30)=30/\varphi(30)=3.75\) reproduce exactly.)}

\subsection{6. Conclusion}\label{conclusion}

The General Divisor Theorem (Theorem 3) gives an exact, elementary
correction to the naive independent-density heuristic for counting
integers in a fixed residue class that are also coprime to a modulus ---
for arbitrary positive integers \(N,m\) and any integer \(a\), with no
divisibility or squarefreeness restriction on either. Writing
\(d=\mathrm{rad}(\gcd(m,N))\) for the primes \(m\) and \(N\) share, and
\(R=\mathrm{rad}(N)/d\) for what remains, the indicator has exact
minimal period \(T_{\min}=\mathrm{lcm}(m,\mathrm{rad}(N))\), with
exactly \(\varphi(R)\) accepted values per period, whenever
\(\gcd(a,d)=1\); the correction factor against the naive prediction is
then exactly \[
C(N,m) = \frac{d}{\varphi(d)},
\] independent of which admissible \(a\) is chosen. If instead
\(\gcd(a,d)>1\), the accepted set is identically empty for every \(L\)
--- a genuine dichotomy, not an edge case. The theorem follows from the
Chinese Remainder Theorem and multiplicativity of \(\varphi\), requiring
no sieve-theoretic or analytic machinery, and its proof was
independently audited and confirmed against exhaustive bounded
computational tests, including the minimal-period claim; verification
artifacts --- along with the audit history --- are available in the
repository (Appendix A).

The primorial case that motivated this work is recovered exactly as one
specialization: for \(N=P_k\) squarefree and \(m\) a squarefree divisor
of \(P_k\), squarefreeness forces \(d=m\), giving \[
C(m) = \frac{m}{\varphi(m)}.
\] The primorial theorem is recovered without modification to its
period, count, or correction factor. In its development, the \(m=6\)
case also corrected an earlier reported constant \(24/25\) to
\(C(6)=3\); against the refined predicted count, the ratio is exactly
\(1\).

Section 4 gives a separate worked example of why a reading should be
stated together with its producing specification. A weighted-vector
construction labeled as corresponding to a \(45^\circ\) diagonal direction
yields a different reading under the stated construction. Together with
the density-correction example, this distinguishes a shared label, a
shared reading, and a shared specification.

\subsection{Appendix A. Verification
Code}\label{appendix-a.-verification-code}

\begin{Shaded}
\begin{Highlighting}[]
\ImportTok{from}\NormalTok{ math }\ImportTok{import}\NormalTok{ gcd}
\ImportTok{from}\NormalTok{ sympy }\ImportTok{import}\NormalTok{ totient, primefactors}

\KeywordTok{def}\NormalTok{ naive\_predicted(limit, primorial, m):}
\NormalTok{    density }\OperatorTok{=}\NormalTok{ totient(primorial) }\OperatorTok{/}\NormalTok{ primorial}
    \ControlFlowTok{return}\NormalTok{ density }\OperatorTok{*}\NormalTok{ (limit }\OperatorTok{/}\NormalTok{ m)}

\KeywordTok{def}\NormalTok{ count\_residue(limit, primorial, m, a):}
\NormalTok{    c }\OperatorTok{=} \DecValTok{0}
\NormalTok{    n }\OperatorTok{=}\NormalTok{ a}
    \ControlFlowTok{while}\NormalTok{ n }\OperatorTok{\textless{}=} \DecValTok{0}\NormalTok{:}
\NormalTok{        n }\OperatorTok{+=}\NormalTok{ m}
    \ControlFlowTok{for}\NormalTok{ n }\KeywordTok{in} \BuiltInTok{range}\NormalTok{(n, limit }\OperatorTok{+} \DecValTok{1}\NormalTok{, m):}
        \ControlFlowTok{if}\NormalTok{ gcd(n, primorial) }\OperatorTok{==} \DecValTok{1}\NormalTok{:}
\NormalTok{            c }\OperatorTok{+=} \DecValTok{1}
    \ControlFlowTok{return}\NormalTok{ c}

\KeywordTok{def}\NormalTok{ refined\_predicted(limit, primorial, m):}
    \CommentTok{\# density = product over primes dividing Pk but NOT dividing m}
\NormalTok{    m\_primes }\OperatorTok{=} \BuiltInTok{set}\NormalTok{(primefactors(m))}
\NormalTok{    dens }\OperatorTok{=} \FloatTok{1.0}
    \ControlFlowTok{for}\NormalTok{ p }\KeywordTok{in}\NormalTok{ primefactors(primorial):}
        \ControlFlowTok{if}\NormalTok{ p }\KeywordTok{not} \KeywordTok{in}\NormalTok{ m\_primes:}
\NormalTok{            dens }\OperatorTok{*=}\NormalTok{ (}\DecValTok{1} \OperatorTok{{-}} \FloatTok{1.0} \OperatorTok{/}\NormalTok{ p)}
    \ControlFlowTok{return}\NormalTok{ dens }\OperatorTok{*}\NormalTok{ (limit }\OperatorTok{/}\NormalTok{ m)}
\end{Highlighting}
\end{Shaded}

Full runnable versions, including the general \((m,a)\) sweep and the
permutation search for Section 4, are in
\texttt{tests/basis/test\_persistent\_constant.py} and
\texttt{tests/basis/test\_9423\_phase\_lock.py}.

\emph{(Added in v3.2.)} The corresponding exhaustive verification for
the general theorem of Section 2 --- covering arbitrary \((N,m,a)\),
both branches of the dichotomy, arbitrary-window invariance, and
divisor-based exact minimal-period detection over more than
\(4.7\times10^6\) checks, with zero counterexamples found --- is in
\texttt{test\_general\_divisor\_theorem.py}, with results recorded in
\texttt{verification-results.md}. The proof itself was independently
audited line by line against a frozen specification
(\texttt{theorem-specification.md}); a first pass
(\texttt{proof-audit.md}) found two genuine gaps --- an unproved
arbitrary-window generalization and an unproved minimality premise ---
which were then closed and confirmed by a second, independent pass
(\texttt{proof-audit-repaired.md}). Neither gap affected the primorial
case (Section 3.4), which was independently reconfirmed as an exact
specialization in \texttt{specialization-audit-repaired.md}.

\subsection{Appendix B. References}\label{appendix-b.-references}

{[}1{]} G. H. Hardy and E. M. Wright, \emph{An Introduction to the
Theory of Numbers}, Oxford University Press. Relevant background: the
Chinese Remainder Theorem, Euler's totient function, and
multiplicativity of \(\varphi\).

\subsection{Appendix C. Project
Vocabulary}\label{appendix-c.-project-vocabulary}

The notation in this appendix is project-specific and is not part of
Theorem 1. The theorem itself remains legible and valid independently of
this vocabulary.

\textbf{Pi --- expand.} \(m \mapsto \mathbb Z/m\mathbb Z\). This exposes
the complete residue/state space for a given modulus before any
restriction is imposed.

\textbf{\(\pi\) --- extend.} \(P_k \mapsto P_{k+1}=p_{k+1}P_k\). This
adds another specification/constraint by extending the primorial by one
further prime.

\textbf{\(\Pi\) --- resist.} \(\prod_{p\mid P_k/m}(1-1/p)\). This is the
fraction of candidates remaining admissible under the divisibility
constraints not already imposed by the residue condition; it is the
quantity appearing directly in the proof of Theorem 1.

These roles give the progression \[
\text{Pi (expand)} \to \pi\text{ (extend)} \to \Pi\text{ (resist)} \to r\text{ (read)}.
\] These names label roles already present in the proof of Theorem 1.
The proof itself does not depend on this vocabulary and holds
independently of it.

\emph{(Added in v3.2.)} This vocabulary describes Theorem 1's proof
specifically. The ``extend'' step (\(P_k\to P_{k+1}\)) does not have a
natural analog in Theorem 3: the general theorem's \(N\) has no
canonical ``next \(N\)'' the way the primorial sequence does, so this
vocabulary should not be read as describing the general theorem of
Section 2.

\begin{center}\rule{0.5\linewidth}{0.5pt}\end{center}

\subsection{\texorpdfstring{Change Log (v3.1 \(\to\) v3.2)}{Change Log (v3.1 -> v3.2)}}\label{change-log-v3.1-v3.2}

Keyed to \texttt{v3.2-integration-audit.md}. Rows below list only
sections where content changed; sections not listed are byte-for-byte
identical to Draft v3.1 (modulo renumbering).

\begin{longtable}[]{@{}
  >{\raggedright\arraybackslash}p{(\columnwidth - 4\tabcolsep) * \real{0.3333}}
  >{\raggedright\arraybackslash}p{(\columnwidth - 4\tabcolsep) * \real{0.3333}}
  >{\raggedright\arraybackslash}p{(\columnwidth - 4\tabcolsep) * \real{0.3333}}@{}}
\toprule\noalign{}
\begin{minipage}[b]{\linewidth}\raggedright
v3.1 location
\end{minipage} & \begin{minipage}[b]{\linewidth}\raggedright
Change made
\end{minipage} & \begin{minipage}[b]{\linewidth}\raggedright
Audit classification
\end{minipage} \\
\midrule\noalign{}
\endhead
\bottomrule\noalign{}
\endlastfoot
Abstract & Added paragraph on the general theorem and its recovery of
Theorem 1 & NEW MATERIAL \\
Section 1 & Added remark: \(d=m\) in the primorial setting, so
\(\gcd(a,d)=1\iff\gcd(a,m)=1\) & WORDING \\
Section 2, after Theorem 1's proof & Added remark pointing to Section 3
and explaining why no minimal-period subtlety arises here & NEW
MATERIAL \\
Section 2, Corollary 1 & Added parenthetical connecting
\(C(6)=m/\varphi(m)\) to \(d/\varphi(d)\) notation & WORDING \\
Section 2, Corollary 2 & Unchanged & NO CHANGE \\
\textbf{New Section 3} & Full new section: Lemma, Theorem 3, Theorem 2,
Theorem~\(1'\), two worked examples, recovery of Theorem 1 & NEW MATERIAL
(largest item) \\
Old Section 3 \(\to\) new Section 4 & Renumbered only; text unchanged & NO
CHANGE \\
Old Section 4 \(\to\) new Section 5, Table 1 & Renumbered; values unchanged;
rechecked against Section 3 formulas, no discrepancy, added confirmation
parenthetical & RECHECK NUMBERS \(\to\) NO CHANGE \\
Old Section 4 \(\to\) new Section 5, \(C(6),C(10),C(30)\) & Renumbered; values
unchanged; rechecked, no discrepancy & RECHECK NUMBERS \(\to\) NO CHANGE \\
Old Section 5 \(\to\) new Section 6 (Conclusion) & Extended with paragraph on
the general theorem, its audit status, and exact recovery of Theorem 1 &
NEW MATERIAL \\
Appendix A & Added pointer to
\texttt{test\_general\_divisor\_theorem.py} & NEW MATERIAL (pointer
only) \\
Appendix B & Unchanged & NO CHANGE \\
Appendix C & Added scope-limiting sentence about the ``extend'' step &
WORDING \\
\end{longtable}

No row involved a \texttt{MATHEMATICAL\ UPDATE} --- consistent with the
integration audit's finding that no existing claim in Draft v3.1
required correction.

\textbf{Note (result-order restructuring):} this Change Log documents
the v3.1\(\to\)v3.2 \emph{content} changes only, as it did before the
result-order restructuring. It describes the section numbers and
positions that existed at the time those changes were made (discovery
order); it has been left unedited, per instruction, rather than updated
to reflect the current physical section numbers. For the record of the
subsequent reordering pass itself (moved blocks, rewritten bridges,
deleted redundancies, and the full cross-reference audit), see
\texttt{v3.2-result-order-audit.md}.

\end{document}
